\chapter{Small Oscillations}

\textbooknumber{\S 21}
\section{Free oscillations in one dimension}

% ============ Problem 1 ============ %

\begin{problem}
{
Express the amplitude and initial phase of the oscillations in terms of the initial co-ordinate $x_0$ and velocity $v_0$.
}
{}
{}
\end{problem}

% ============ Problem 2 ============ %

\begin{problem}
{
Find the ratio of frequencies $\omega$ and $\omega'$ of the oscillations of two diatomic molecules consisting of atoms of different isotopes, the masses of the atoms being $m_1$, $m_2$ and $m_1'$, $m_2'$.
}
{}
{}
\end{problem}

% ============ Problem 3 ============ %

\begin{problem}
{
Find the frequency of oscillations of a particle of mass $m$ which is free to move along a line and is attached to a spring whose other end is fixed at a point $A$ (Fig.~22) at a distance $l$ from the line. A force $F$ is required to extend the spring to length $l$.
}
{}
{}
\end{problem}

% ============ Problem 4 ============ %

\begin{problem}
{
The same as \probref{\S 21}{3}, but for a particle of mass $m$ moving on a circle of radius $r$ (Fig.~23).
}
{}
{}
\end{problem}

% ============ Problem 5 ============ %

\begin{problem}
{
Find the frequency of oscillations of the pendulum shown in the figure in \probref{\S 5}{2}, whose point of support carries a mass $m_1$ and is free to move horizontally (its motion is integrated in \probref{\S 14}{3}).
}
{}
{}
\end{problem}

% ============ Problem 6 ============ %

\begin{problem}
{
Determine the form of a curve such that the frequency of oscillations of a particle on it under the force of gravity is independent of the amplitude.
}
{}
{}
\end{problem}

\textbooknumber{\S 22}
\section{Forced oscillations}

% ============ Problem 1 ============ %

\begin{problem}
{
Determine the forced oscillations of a system under a force $F(t)$ of the following forms, if at time $t = 0$ the system is at rest in equilibrium ($x = \mathrm{d}x/\mathrm{d}t = 0$):
}{}{}
\end{problem}
\begin{subproblem}
{
$F = F_0$, a constant.
}
{}
{}
\end{subproblem}

\begin{subproblem}
{
$F = at$.
}
{}
{}
\end{subproblem}

\begin{subproblem}
{
$F = F_0 e^{-\alpha t}$.
}
{}
{}
\end{subproblem}

\begin{subproblem}
{
$F = F_0 e^{-\alpha t} \cos{\beta t}$.
}
{}
{}
\end{subproblem}

% ============ Problem 2 ============ %

\begin{problem}
{
Determine the final amplitude for the oscillations of a system under a force which is zero for $t < 0$, $F_0 t/T$ for $0 < t < T$, and $F_0$ for $t > T$ (Fig.~24), if up to time $t = 0$ the system is at rest in equilibrium.
}
{}
{}
\end{problem}

% ============ Problem 3 ============ %

\begin{problem}
{
The same as \probref{\S 22}{2}, but for a constant force $F_0$ which acts for a finite time $T$ (Fig.~25).
}
{}
{}
\end{problem}

% ============ Problem 4 ============ %

\begin{problem}
{
The same as \probref{\S 22}{2}, but for a force $F_0 t/T$ which acts between $t = 0$ and $t = T$ (Fig.~26).
}
{}
{}
\end{problem}

% ============ Problem 5 ============ %

\begin{problem}
{
The same as \probref{\S 22}{2}, but for a force $F_0 \sin{\omega t}$ which acts between $t = 0$ and $t = T = 2\pi/\omega$ (Fig.~27).
}
{}
{}
\end{problem}

\textbooknumber{\S 23}
\section{Oscillations of systems with more than one degree of freedom}

% ============ Problem 1 ============ %

\begin{problem}
{
Determine the oscillations of a system with two degrees of freedom whose Lagrangian is
\begin{equation*}
    L = \frac{1}{2}\left(\Dot{x}^2 + \Dot{y}^2\right) - \frac{1}{2}\omega_0^2\left(x^2 + y^2\right) + \alpha x y
\end{equation*}
(two identical one-dimensional systems of eigenfrequency $\omega_0$ coupled by an interaction $-\alpha x y$).
}
{}
{}
\end{problem}

% ============ Problem 2 ============ %

\begin{problem}
{
Determine the small oscillations of a coplanar double pendulum (see the figure in \probref{\S 5}{1}).
}
{}
{}
\end{problem}

% ============ Problem 3 ============ %

\begin{problem}
{
Find the path of a particle in a central field $U = \frac{1}{2} k r^2$ (called a space oscillator).
}
{}
{}
\end{problem}

\textbooknumber{\S 24}
\section{Vibrations of molecules}

% ============ Problem 1 ============ %

\begin{problem}
{
Determine the frequencies of vibrations of a symmetrical linear triatomic molecule $ABA$ (Fig.~28). It is assumed that the potential energy of the molecule depends only on the distances $AB$ and $BA$ and the angle $ABA$.
}
{}
{}
\end{problem}

% ============ Problem 2 ============ %

\begin{problem}
{
The same as \probref{\S 24}{1}, but for a triangular molecule $ABA$ (Fig.~29).
}
{}
{}
\end{problem}

% ============ Problem 3 ============ %

\begin{problem}
{
The same as \probref{\S 24}{1}, but for an unsymmetrical linear molecule $ABC$ (Fig.~30).
}
{}
{}
\end{problem}

\textbooknumber{\S 26}
\section{Forced oscillations under friction}

% ============ Problem 1 ============ %

\begin{problem}
{
Determine the forced oscillations due to an external force $f = f_0 e^{\alpha t} \cos{\gamma t}$ in the presence of friction.
}
{}
{}
\end{problem}

\textbooknumber{\S 27}
\section{Parametric resonance}

% ============ Problem 1 ============ %

\begin{problem}
{
Obtain an expression correct as far as the term in $h^2$ for the limits of the region of instability for resonance near $\gamma = 2\omega_0$.
}
{}
{}
\end{problem}

% ============ Problem 2 ============ %

\begin{problem}
{
Determine the limits of the region of instability in resonance near $\gamma = \omega_0$.
}
{}
{}
\end{problem}

% ============ Problem 3 ============ %

\begin{problem}
{
Find the conditions for parametric resonance in small oscillations of a simple pendulum whose point of support oscillates vertically (see the Lagrangian in \probref{\S 5}{3}(c)).
}
{}
{}
\end{problem}

\textbooknumber{\S 29}
\section{Resonance in non-linear oscillations}

% ============ Problem 1 ============ %

\begin{problem}
{
Determine the function $b(\epsilon)$ for resonance at frequencies $\gamma \approx 3\omega_0$.
}
{}
{}
\end{problem}

\textbooknumber{\S 30}
\section{Motion in a rapidly oscillating field}

% ============ Problem 1 ============ %

\begin{problem}
{
Determine the positions of stable equilibrium of a pendulum whose point of support oscillates vertically with a high frequency $\gamma \gg \sqrt{g/l}$ (see the Lagrangian in \probref{\S 5}{3}(c)).
}
{}
{}
\end{problem}

% ============ Problem 2 ============ %

\begin{problem}
{
The same as \probref{\S 30}{1}, but for a pendulum whose point of support oscillates horizontally (see the Lagrangian in \probref{\S 5}{3}(b)).
}
{}
{}
\end{problem}
